\documentclass[11pt]{article}
\usepackage[review]{acl}
\usepackage{times}
\usepackage{latexsym}
\usepackage[T1]{fontenc}
\usepackage[utf8]{inputenc}
\usepackage{microtype}
\usepackage{inconsolata}
\usepackage{graphicx}
\usepackage{booktabs}
\usepackage{amsmath}
\usepackage[indonesian]{babel}
\usepackage{xurl}

\title{StereoQueerEval: Klasifikasi Ujaran Kebencian terhadap\
Komunitas LGBTQIA+ dengan XLM-RoBERTa dan Qwen3.5-9B}
% Identitas dan afiliasi diisi sebelum beralih ke mode final.
\author{Penulis Anonim}

\begin{document}
% Nomor halaman dipertahankan; nomor baris review dinonaktifkan.
\nolinenumbers
\raggedbottom
\maketitle
\begin{abstract}
% Evidence: Q001,Q002,Q004,Q005,B001.
Penelitian ini menyajikan pendekatan untuk menyelesaikan Subtask B SemEval-2027 Task 8: StereoQueerEval, yaitu mengklasifikasikan komentar YouTube ke dalam kelas \texttt{no}, \texttt{yes\_implicit}, dan \texttt{yes\_explicit}. XLM-RoBERTa Base digunakan sebagai pembanding, sedangkan pendekatan utama menyesuaikan Qwen3.5-9B menggunakan Low-Rank Adaptation (LoRA). Komentar dipertahankan dalam bahasa aslinya dan diproses tanpa judul serta deskripsi video. Qwen setelah pelatihan berlabel dievaluasi menggunakan prompt Inggris, Italia, dan Belanda. Seluruh pilihan prompt Qwen menghasilkan macro-F1 lebih tinggi daripada XLM-RoBERTa pada pengujian internal yang sama. Nilai macro-F1 tertinggi yang teramati adalah 0,719146 untuk Inggris, 0,7036 untuk Italia, dan 0,641522 untuk Belanda. Pada hasil per kelas yang tersedia, \texttt{no} memperoleh F1 tertinggi dan \texttt{yes\_implicit} memperoleh F1 terendah dalam setiap pasangan bahasa data dan prompt. Kesamaan bahasa prompt dan komentar belum menghasilkan skor tertinggi secara konsisten. Hasil ini berasal dari pengujian internal dan belum merupakan hasil evaluasi resmi shared task.
\end{abstract}

\section{Pendahuluan}
\label{sec:pendahuluan}
% Evidence: L102.
Deteksi ujaran kebencian tidak hanya berkaitan dengan pengenalan kata-kata kasar. Pesan yang menyerang kelompok tertentu juga dapat disampaikan melalui ungkapan tersirat yang membutuhkan pemahaman konteks. \citet{elsherief-etal-2021-latent} menjelaskan bahwa bahasa berkode dan ungkapan tidak langsung menjadi tantangan dalam mendeteksi ujaran kebencian yang disampaikan secara tidak langsung. Perbedaan cara penyampaian tersebut membuat model perlu mengenali maksud pesan, bukan hanya kata-kata yang muncul di dalamnya.

% Evidence: L111.
\citet{chan-etal-2024-hate} mengkaji deteksi ujaran kebencian terhadap komunitas LGBTQIA+ menggunakan GPT-3.5-turbo pada teks asli dan teks yang diterjemahkan ke bahasa Inggris. Penerjemahan memberikan hasil yang beragam; sebagian istilah slang dan ungkapan yang berkaitan dengan budaya tertentu tidak diterjemahkan dengan tepat. Temuan tersebut menjadi pertimbangan untuk mempertahankan bahasa asli komentar. Penelitian itu juga menyarankan penggunaan prompt dalam bahasa data serta pengujian model lain. Saran tersebut menjadi salah satu dasar penyusunan prompt klasifikasi dalam penelitian ini.

% Evidence: L110,L117,Q001,B001.
Penelitian ini berfokus pada penyelesaian \textit{Subtask B: Hate Speech Classification} dalam SemEval-2027 Task 8: StereoQueerEval. Teks yang diklasifikasikan berupa komentar berbahasa Inggris, Italia, dan Belanda. Setiap komentar diberi label \texttt{no}, \texttt{yes\_implicit}, atau \texttt{yes\_explicit}. XLM-RoBERTa Base digunakan sebagai pembanding, sedangkan pendekatan utama menggunakan Qwen3.5-9B yang disesuaikan dengan data task. Keduanya dinilai pada komentar uji yang sama.

% Evidence: Q001,Q002,Q004,Q006,B001.
Penyelesaian task dilakukan melalui penyiapan data, pelatihan model, dan penentuan satu label untuk setiap komentar. Pada Qwen, definisi kategori diberikan melalui prompt berbahasa Inggris, Italia, atau Belanda. Ketiga pilihan digunakan untuk menemukan pengaturan yang sesuai bagi klasifikasi, dengan tetap mempertahankan bahasa asli komentar. Hasil evaluasi dilaporkan sebagai macro-F1 pada setiap bahasa data dan F1 setiap kelas, sesuai nilai yang tersedia dalam ringkasan eksperimen.

\section{Deskripsi Task}
\subsection{SemEval-2027 Task 8: StereoQueerEval}
\label{sec:task}
% Evidence: L110,L117.
StereoQueerEval merupakan task identifikasi stereotip dan ujaran kebencian yang menargetkan komunitas LGBTQIA+.\footnote{\url{https://semeval.github.io/SemEval2027/tasks.html}} Data berupa komentar YouTube yang disertai judul dan deskripsi video sebagai informasi konteks. Task ini terdiri atas deteksi stereotip, klasifikasi ujaran kebencian, serta identifikasi target dan identitas LGBTQIA+ yang menjadi sasaran kebencian \citep{stereoqueereval-2026-data}.

Penelitian ini berfokus pada Subtask B. Setiap komentar diberi salah satu dari tiga label berikut:
\begin{itemize}
    \item \texttt{no}: komentar yang tidak mengandung ujaran kebencian, termasuk komentar yang mendukung komunitas LGBTQIA+.
    \item \texttt{yes\_implicit}: kebencian yang disampaikan secara tersirat atau tidak langsung.
    \item \texttt{yes\_explicit}: kebencian yang dinyatakan secara langsung terhadap individu atau kelompok LGBTQIA+.
\end{itemize}

% Evidence: L110,Q001,Q005.
Data StereoQueerEval tersedia dalam bahasa Inggris, Italia, Belanda, dan Persia. Percobaan ini menggunakan tiga bahasa pertama, yang selanjutnya disingkat EN, IT, dan NL. Masukan dibatasi pada teks komentar beserta penanda bahasa; judul, deskripsi video, dan anotasi dari subtask lain tidak digunakan untuk menentukan label.

\subsection{Penelitian Terkait}
% Evidence: L101,L114,L112.
Deteksi kebencian terhadap komunitas LGBTQIA+ telah dikaji dengan model klasifikasi maupun model bahasa besar. \citet{nozza-2022-nozza} menggabungkan prediksi BERT, RoBERTa, dan HateBERT untuk mengenali homofobia dan transfobia pada komentar YouTube Inggris. QUEEREOTYPES menyediakan anotasi stereotip dan ujaran kebencian secara terpisah pada data Italia \citep{cignarella-etal-2024-queereotypes}. Sementara itu, \citet{nguyen-etal-2025-leveraging} mengevaluasi lima model bahasa besar untuk mengenali teks anti-LGBTQIA+ tanpa pelatihan tambahan dan melaporkan bahwa kesalahan klasifikasi masih terjadi.

% Evidence: L111,L115,L104.
\citet{chan-etal-2024-hate} membandingkan GPT-3.5-turbo sebelum dan sesudah pelatihan pada data berlabel. Hasilnya mendukung penggunaan pelatihan untuk klasifikasi multibahasa, tetapi tugasnya menggunakan dua kelas dan weighted-F1. Penelitian ini membedakan tiga kelas dan menggunakan macro-F1, sehingga nilai kedua penelitian tidak dibandingkan secara langsung. XLM-RoBERTa dipilih sebagai pembanding karena prapelatihannya menggunakan teks dari 100 bahasa \citep{conneau-etal-2020-unsupervised}, sedangkan LoRA memungkinkan penyesuaian Qwen melalui sejumlah kecil parameter tambahan \citep{hu-etal-2021-lora}.

% Evidence: L116,Q001,Q006.
Pilihan bahasa prompt juga perlu diperhatikan ketika menggunakan model multibahasa. \citet{behzad-etal-2024-ask} menemukan bahwa prompt dalam bahasa berbeda dapat menghasilkan nilai yang berbeda pada penilaian tata bahasa Inggris. Temuan tersebut berasal dari tugas lain, sehingga menjadi pertimbangan untuk menguji pilihan prompt, bukan bukti langsung tentang deteksi kebencian. Dalam penyelesaian StereoQueerEval ini, tiga bahasa prompt digunakan pada tahap pelatihan dan pengujian Qwen.

\section{Metodologi}
\subsection{XLM-RoBERTa}
% Evidence: L115,B001.
XLM-RoBERTa merupakan model Transformer multibahasa yang dilatih menggunakan lebih dari 2 TB teks CommonCrawl dari 100 bahasa \citep{conneau-etal-2020-unsupervised}. Varian Base digunakan sebagai pembanding dalam penelitian ini. Teks komentar diubah menjadi token, kemudian diproses untuk memperoleh representasi yang mempertimbangkan hubungan antarkata. Representasi tersebut diteruskan ke lapisan klasifikasi untuk memprediksi tiga kelas task.

Seluruh parameter XLM-RoBERTa dan lapisan klasifikasi diperbarui menggunakan data pelatihan. Pelatihan menggunakan bobot kelas yang dihitung dari jumlah contoh pada data pelatihan agar kelas yang lebih sedikit mendapat perhatian dalam perhitungan kesalahan. Model dengan macro-F1 validasi tertinggi digunakan untuk pengujian.

\subsection{Qwen3.5-9B}
% Evidence: Q001,Q004,L103,L104.
Qwen3.5-9B digunakan sebagai model bahasa yang menerima prompt klasifikasi \citep{qwen-team-2026-qwen35}. Model disesuaikan melalui \textit{supervised fine-tuning} (SFT), yaitu pelatihan menggunakan komentar dan jawaban yang benar. Penyesuaian dilakukan dengan LoRA, yang mempertahankan bobot utama model dan memperbarui parameter tambahan berukuran kecil \citep{hu-etal-2021-lora}. Sekitar 0,48\% dari parameter yang dimuat diperbarui selama pelatihan.

% Evidence: Q001,Q004,Q006.
Prompt menjelaskan ketiga kategori dan menegaskan bahwa penyebutan identitas LGBTQIA+ saja tidak cukup untuk menganggap komentar sebagai kebencian. Prompt tersedia dalam bahasa Inggris, Italia, dan Belanda, tanpa contoh komentar berlabel. Jawaban yang diminta berupa satu angka: 0 untuk \texttt{no}, 1 untuk \texttt{yes\_explicit}, dan 2 untuk \texttt{yes\_implicit}. Teks komentar tetap menggunakan bahasa aslinya. Penanda bahasa dan komentar pada masukan tetap ditulis dalam bahasa Inggris untuk seluruh pilihan prompt.

% Evidence: Q001,Q004.
Pelatihan diarahkan pada jawaban yang benar, sedangkan prompt dan komentar tidak dimasukkan ke perhitungan kesalahan pelatihan. Saat prediksi, model menghitung jumlah log probabilitas setiap pilihan jawaban beserta penanda akhirnya. Jawaban dengan skor tertinggi menjadi label prediksi. Dengan aturan ini, hasil selalu berupa salah satu dari tiga kategori dan tidak bergantung pada penjelasan bebas dari model.

\subsection{Alur Analisis}
% Evidence: Q001,Q004,Q005,Q006,B001.
Alur analisis meliputi eksplorasi data, prapemrosesan dan pembagian data, pelatihan dan prediksi, serta evaluasi model. XLM-RoBERTa mempelajari klasifikasi langsung dari teks komentar, sedangkan Qwen menerima definisi kategori dan komentar dalam format percakapan. Kedua pendekatan menghasilkan satu label untuk setiap komentar. Langkah analisis diterapkan pada setiap bahasa data secara terpisah.

\subsubsection{Eksplorasi Data (EDA)}
% Evidence: Q001,Q005,B001.
Eksplorasi data atau exploratory data analysis (EDA) dilakukan untuk memeriksa kelengkapan komentar, keunikan ID, kesesuaian label, dan jumlah contoh setiap kelas pada setiap bahasa. Pemeriksaan juga dilakukan untuk mengetahui apakah komentar atau konteks video yang sama muncul pada bagian data berbeda. Hasil pemeriksaan digunakan untuk memahami komposisi data dan memastikan pembagian yang digunakan sesuai untuk pelatihan serta evaluasi.

\subsubsection{Prapemrosesan dan Pembagian Data}
% Evidence: Q001,Q004,Q005,B001.
Praproses mengikuti pengaturan masing-masing pendekatan. Pada XLM-RoBERTa, huruf diseragamkan menjadi huruf kecil melalui \textit{casefold} dan emoji diganti dengan penanda \texttt{[EMOJI]}. Pada Qwen, kapitalisasi, emoji, dan tanda baca dipertahankan. Komentar tidak diterjemahkan. Perbedaan praproses ini diperhitungkan ketika menafsirkan perbandingan kedua pendekatan.

Normalisasi Unicode NFKC, casefold, dan perapian spasi juga digunakan untuk memeriksa kesamaan komentar antarbagian data. Normalisasi untuk pemeriksaan tersebut tidak mengubah teks masukan Qwen. Sebelum tokenisasi, kelengkapan komentar, keunikan ID, dan kesesuaian label diperiksa. Masukan XLM-RoBERTa dibatasi hingga 384 token, sedangkan masukan Qwen hingga 1.024 token termasuk prompt dan jawaban. Pada Qwen, komentar yang terlalu panjang dipotong dengan tetap menyediakan ruang untuk prompt dan pilihan jawaban.

% Evidence: Q005,B001.
Data setiap bahasa dibagi menjadi bagian pelatihan, validasi, dan pengujian dengan perbandingan mendekati 70:10:20 dan seed 42. Pembagian mempertahankan kelompok komentar dengan konteks video yang sama serta komentar duplikat pada bagian yang sama. Kelompok konteks dikenali melalui pasangan judul dan deskripsi video. Karena kelompok tidak dipecah, jumlah akhir tidak selalu tepat mengikuti rasio tersebut.

\subsubsection{Pelatihan dan Prediksi}
% Evidence: Q001,Q004,Q006,B001.
Data pelatihan digunakan untuk memperbarui parameter model. Setelah setiap epoch, hasil pada data validasi dihitung untuk memilih hasil pelatihan dengan macro-F1 tertinggi. Model terpilih kemudian memprediksi label komentar uji menggunakan aturan klasifikasi masing-masing model. Label yang benar pada data uji hanya digunakan untuk menghitung metrik evaluasi.

\subsubsection{Evaluasi Model}
% Evidence: Q001,Q003,Q004,Q005,B001.
Evaluasi dilakukan setelah model dipilih berdasarkan hasil validasi. Label prediksi dibandingkan dengan label yang benar pada komentar uji yang sama untuk kedua pendekatan. Hasil disajikan sebagai macro-F1 untuk setiap pasangan bahasa data dan prompt, serta F1 setiap kelas. Nilai yang belum tersedia dalam ringkasan eksperimen dibiarkan kosong. Pengujian ini merupakan evaluasi internal karena komentar uji berasal dari pembagian data pelatihan penyelenggara, bukan dari data uji resmi shared task.

\subsection{Metrik Evaluasi}
% Evidence: Q001,Q002,Q003,Q004,B001.
Ukuran utama adalah macro-F1, yaitu rata-rata F1 dari tiga kelas dengan bobot yang sama. Precision menunjukkan proporsi prediksi suatu kelas yang benar, sedangkan recall menunjukkan proporsi contoh kelas tersebut yang berhasil dikenali. F1 menggabungkan precision dan recall sehingga dapat merangkum keduanya.

\section{Rancangan Eksperimen}
% Evidence: Q005,B001.
Data yang dianalisis berjumlah 2.989 komentar Inggris, 2.400 Italia, dan 2.238 Belanda. Tabel~\ref{tab:data} menyajikan jumlah komentar pada setiap kelas dan bagian data. Jumlah komentar uji adalah 586 untuk Inggris, 481 untuk Italia, dan 456 untuk Belanda. Identitas dan label komentar uji yang digunakan kedua pendekatan telah diperiksa dan sama.

\begin{table}[t]
\centering
\small
\setlength{\tabcolsep}{3pt}
\begin{tabular}{llrrr}
\toprule
Bahasa & Kelas & train & val & test \\
\midrule
EN & \texttt{no} & 983 & 147 & 277 \\
EN & \texttt{yes\_implicit} & 546 & 70 & 156 \\
EN & \texttt{yes\_explicit} & 582 & 75 & 153 \\
EN & Total & 2111 & 292 & 586 \\
IT & \texttt{no} & 822 & 105 & 233 \\
IT & \texttt{yes\_implicit} & 401 & 56 & 124 \\
IT & \texttt{yes\_explicit} & 476 & 59 & 124 \\
IT & Total & 1699 & 220 & 481 \\
NL & \texttt{no} & 784 & 109 & 234 \\
NL & \texttt{yes\_implicit} & 335 & 55 & 97 \\
NL & \texttt{yes\_explicit} & 437 & 62 & 125 \\
NL & Total & 1556 & 226 & 456 \\
\bottomrule

\end{tabular}
\caption{Jumlah komentar pada pembagian data menurut label kelas. train, val, dan test berarti pelatihan, validasi, dan pengujian.}
\label{tab:data}
\end{table}


% Evidence: Q001,Q002,Q004,Q006,B001.
Eksperimen membandingkan XLM-RoBERTa Base dengan Qwen3.5-9B yang disesuaikan menggunakan SFT LoRA. XLM-RoBERTa dilatih satu kali untuk setiap bahasa, sedangkan Qwen dilatih untuk sembilan pasangan bahasa komentar dan bahasa prompt. Setiap pasangan Qwen dimulai dari model dasar yang sama dengan parameter LoRA baru. Prompt yang dipilih digunakan secara konsisten selama pelatihan, validasi, dan pengujian. Pemisahan data mengikuti prosedur pada bagian Metodologi.

% Evidence: Q001,Q004,B001,L109.
Kedua pendekatan menggunakan AdamW \citep{loshchilov-hutter-2019-decoupled}, seed 42, dan weight decay sebesar 0,01. XLM-RoBERTa dilatih selama paling banyak 10 epoch dengan learning rate $10^{-5}$, batch 16, dan batas masukan 384 token. Learning rate diturunkan secara linear tanpa tahap pemanasan. Pelatihan menggunakan presisi FP16 pada NVIDIA GeForce RTX 3060 Ti.

% Evidence: Q001,Q004.
Qwen dilatih selama paling banyak lima epoch dengan learning rate $10^{-4}$ dan batas masukan 1.024 token. Setiap batch berisi empat komentar dengan akumulasi empat langkah, sehingga batch efektifnya adalah 16. LoRA menggunakan rank 16, alpha 32, dan dropout 0,05. Learning rate mengikuti penjadwalan cosine dengan tahap pemanasan sekitar 5\% dari rencana langkah pelatihan. Perhitungan menggunakan BF16 pada NVIDIA H200.

% Evidence: Q001,Q002,Q003,Q004,B001.
Evaluasi dilakukan pada akhir setiap epoch. Pelatihan dihentikan lebih awal jika macro-F1 validasi tidak membaik dalam tiga evaluasi untuk XLM-RoBERTa atau dua evaluasi untuk Qwen. Model dengan macro-F1 validasi tertinggi digunakan untuk pengujian. Pelaporan hasil mengikuti nilai macro-F1 dan F1 per kelas yang tersedia dalam ringkasan eksperimen. Nilai tertinggi pada pengujian dilaporkan secara deskriptif, sedangkan pemilihan pengaturan untuk pengujian berikutnya perlu menggunakan data validasi.

\section{Hasil dan Pembahasan}
\subsection{Hasil Klasifikasi}
% Evidence: Q002,B001.
Tabel~\ref{tab:main} menyajikan macro-F1 XLM-RoBERTa dan Qwen3.5-9B setelah SFT pada sembilan pasangan bahasa data dan prompt. Pada data Inggris, hasil tertinggi diperoleh dengan prompt Inggris sebesar 0,719146. Pada Italia, prompt Inggris memperoleh skor tertinggi sebesar 0,7036, sedangkan pada Belanda skor tertinggi diperoleh dengan prompt Belanda sebesar 0,641522. Nilai tersebut menunjukkan hasil tertinggi yang teramati pada pengujian internal, bukan pengaturan akhir yang dipilih melalui data validasi.

\begin{table*}[t]
\centering
\small
\setlength{\tabcolsep}{3pt}
\begin{tabular}{lrrrr}
\toprule
Data & XLM-RoBERTa & Prompt EN & Prompt IT & Prompt NL \\
\midrule
EN & 0.484506 & \textbf{0.719146} & 0.715757 & 0.688340 \\
IT & 0.527877 & \textbf{0.7036} & 0.702038 & 0.701449 \\
NL & 0.512953 & 0.629927 & 0.630810 & \textbf{0.641522} \\
\bottomrule

\end{tabular}
\caption{Macro-F1 XLM-RoBERTa sebagai pembanding dan Qwen setelah SFT. Nilai Qwen mengikuti ringkasan eksperimen. Cetak tebal menunjukkan nilai Qwen tertinggi pada setiap bahasa data.}
\label{tab:main}
\end{table*}

% Evidence: B001,Q002.
XLM-RoBERTa memperoleh macro-F1 sebesar 0,484506 pada Inggris, 0,527877 pada Italia, dan 0,512953 pada Belanda. Semua pilihan prompt Qwen menghasilkan skor lebih tinggi daripada pembanding tersebut. Hasil kedua pendekatan berasal dari komentar uji yang sama. Namun, ukuran model, praproses, dan cara pelatihannya berbeda, sehingga selisih skor belum dapat dikaitkan dengan satu komponen tertentu.

\subsection{Analisis Berdasarkan Bahasa}
% Evidence: Q002.
Macro-F1 pada data Inggris berada pada rentang 0,688340--0,719146. Selisih hasil tertinggi dan terendah sekitar 0,0308. Prompt Inggris memberikan hasil tertinggi, diikuti prompt Italia sebesar 0,715757 dan prompt Belanda sebesar 0,688340.

Pada data Italia, ketiga prompt menghasilkan skor yang berdekatan, yaitu 0,701449--0,7036, dengan selisih sekitar 0,0022. Prompt Inggris memperoleh skor tertinggi, diikuti prompt Italia sebesar 0,702038 dan prompt Belanda sebesar 0,701449. Dengan demikian, prompt yang menggunakan bahasa yang sama dengan komentar belum menghasilkan skor tertinggi pada data ini.

Pada data Belanda, macro-F1 berada pada rentang 0,629927--0,641522. Prompt Belanda memberikan hasil tertinggi, diikuti prompt Italia sebesar 0,630810 dan prompt Inggris sebesar 0,629927. Selisih hasil tertinggi dan terendah sekitar 0,0116.

Ketiga bahasa memperlihatkan pilihan prompt terbaik yang berbeda. Kesamaan bahasa prompt dan komentar menghasilkan skor tertinggi pada Inggris dan Belanda, sedangkan Italia memperoleh skor tertinggi dengan prompt Inggris. Nilai pada data Belanda juga lebih rendah dibandingkan Inggris dan Italia pada seluruh pilihan prompt. Hasil ini bersifat deskriptif dan belum cukup untuk menjelaskan penyebab perbedaan antarbahasa.

\subsection{Analisis Berdasarkan Kelas}
% Evidence: Q002.
Tabel~\ref{tab:classes} menyajikan F1 setiap kelas untuk pasangan bahasa data dan prompt yang tersedia. Kelas \texttt{no} memperoleh F1 tertinggi dalam setiap pasangan yang memiliki hasil lengkap, sedangkan \texttt{yes\_implicit} memperoleh F1 terendah. Nilai F1 mengikuti ringkasan eksperimen. Sel yang belum memiliki nilai tetap dibiarkan kosong.

\begin{table*}[t]
\centering
\begin{tabular}{llrrr}
\toprule
Data & Prompt & \texttt{no} & \texttt{yes\_implicit} & \texttt{yes\_explicit} \\
\midrule
EN & EN & \textbf{0.8406} & \textbf{0.6172} & 0.6996 \\
EN & IT & 0.8229 & 0.5882 & \textbf{0.7362} \\
EN & NL & 0.8250 & 0.5397 & 0.7003 \\
IT & EN & \textbf{0.7876} & \textbf{0.6127} & 0.7104 \\
IT & IT &  &  &  \\
IT & NL & 0.7806 & 0.6020 & \textbf{0.7217} \\
NL & EN & 0.8118 & \textbf{0.5094} & 0.5686 \\
NL & IT & 0.7908 & 0.4806 & 0.6210 \\
NL & NL & \textbf{0.8140} & 0.4828 & \textbf{0.6278} \\
\bottomrule

\end{tabular}
\caption{F1 setiap kelas setelah SFT, dengan empat angka desimal sesuai sumber. Cetak tebal mengikuti nilai tertinggi yang tersedia untuk setiap kelas pada masing-masing bahasa data. Sel kosong berarti nilai belum tersedia.}
\label{tab:classes}
\end{table*}

% Evidence: Q002.
Pada Inggris, prompt Inggris memberikan hasil tertinggi pada dua kelas, yaitu \texttt{no} dan \nolinkurl{yes_implicit}. Sementara itu, skor tertinggi pada kelas \nolinkurl{yes_explicit} diperoleh dengan prompt Italia. Dengan demikian, prompt yang memberikan macro-F1 tertinggi belum tentu memberikan F1 tertinggi pada seluruh kelas.

Pada Italia, hasil per kelas tersedia untuk prompt Inggris dan Belanda. Prompt Inggris memberikan hasil lebih tinggi pada kelas \texttt{no} dan \texttt{yes\_implicit}. Untuk kelas \texttt{yes\_explicit}, prompt Belanda memberikan hasil lebih tinggi. Baris prompt Italia belum memiliki nilai F1 setiap kelas, sehingga perbandingan pada bahasa ini masih terbatas pada dua prompt.

Pada Belanda, prompt Belanda memperoleh hasil tertinggi pada kelas \texttt{no} dan \texttt{yes\_explicit}. Untuk kelas \texttt{yes\_implicit}, skor tertinggi diperoleh dengan prompt Inggris sebesar 0,5094. Hasil ini menunjukkan bahwa prompt dengan macro-F1 tertinggi tidak selalu memberikan F1 tertinggi pada setiap kelas.

\section{Kesimpulan}
% Evidence: Q001,Q002,Q004,B001.
Penyelesaian Subtask B StereoQueerEval dilakukan melalui klasifikasi komentar Inggris, Italia, dan Belanda ke dalam tiga kelas. Pendekatan utama menyesuaikan Qwen3.5-9B menggunakan SFT LoRA, dengan tiga pilihan bahasa prompt. XLM-RoBERTa tetap digunakan sebagai pembanding. Seluruh pilihan prompt Qwen menghasilkan macro-F1 lebih tinggi pada data uji yang sama. Nilai Qwen tertinggi yang teramati adalah 0,719146 pada Inggris, 0,7036 pada Italia, dan 0,641522 pada Belanda.

% Evidence: Q002.
Kesamaan bahasa prompt dan komentar belum memberikan hasil tertinggi secara konsisten. Pada hasil per kelas yang tersedia, \texttt{no} memperoleh F1 tertinggi dan \texttt{yes\_implicit} memperoleh F1 terendah dalam setiap pasangan bahasa data dan prompt. Selain itu, prompt dengan macro-F1 tertinggi belum tentu memberikan hasil tertinggi pada seluruh kelas. Pengembangan berikutnya perlu melengkapi hasil yang belum tersedia, mengulang percobaan dengan beberapa seed, dan memilih pengaturan berdasarkan data validasi sebelum pengujian lanjutan.

\section*{Keterbatasan}
% Evidence: Q001,Q002,Q006.
Hasil Qwen dibahas berdasarkan macro-F1 dan F1 per kelas yang tersedia dalam ringkasan eksperimen. Perbandingan dengan XLM-RoBERTa menggunakan hasil pembanding yang telah diperiksa pada komentar uji yang sama. Hasil per kelas untuk data Italia dengan prompt Italia belum tersedia dan dibiarkan kosong. Oleh karena itu, perbandingan kelas pada data Italia hanya dilakukan untuk prompt Inggris dan Belanda. Nilai F1 saja belum dapat menjelaskan penyebab kesalahan pada masing-masing komentar.

Percobaan menggunakan satu seed dan satu versi prompt per bahasa. Setiap pasangan bahasa memiliki adapter yang dilatih terpisah, sehingga hasil belum menunjukkan perbedaan bahasa prompt pada satu model terlatih yang sama. Data uji berasal dari pembagian internal data pelatihan penyelenggara. Hasil belum mencerminkan kinerja pada data uji resmi maupun platform dan topik lain.

\bibliography{references}
\end{document}
